% Table benchmark fixture for ReadableSTEM (written for this benchmark; no third-party text).
\documentclass[11pt]{article}
\usepackage[OT1]{fontenc}
\renewcommand{\rmdefault}{cmr}
\usepackage{amsmath}
\usepackage{booktabs}
\usepackage{multirow}
\usepackage{siunitx}
\usepackage{array}
\usepackage[margin=1in]{geometry}
\pagestyle{plain}
\begin{document}

\section*{Lab 4: The Simple Pendulum}

Each group timed twenty swings of a pendulum for four string lengths and computed the
gravitational acceleration from the period. The period of a simple pendulum for small
angles is
\begin{equation}
  T = 2\pi\sqrt{\frac{L}{g}}, \qquad g = \frac{4\pi^2 L}{T^2}.
\end{equation}
Table~\ref{tab:pendulum} lists the measured values.

\begin{table}[h]
\centering
\caption{Measured periods and the derived value of $g$.}
\label{tab:pendulum}
\begin{tabular}{c S[table-format=1.3] S[table-format=1.3] S[table-format=2.2]}
\toprule
Trial & {Length (\si{\metre})} & {Period (\si{\second})} & {$g$ (\si{\metre\per\second\squared})} \\
\midrule
1 & 0.250 & 1.004 & 9.79 \\
2 & 0.500 & 1.420 & 9.79 \\
3 & 0.750 & 1.737 & 9.81 \\
4 & 1.000 & 2.006 & 9.81 \\
\bottomrule
\end{tabular}
\end{table}

The values agree within one per cent. The uncertainty in each period comes mostly from
the reaction time of the person holding the stopwatch, which is about \SI{0.1}{\second}
for a single swing but is divided by twenty when twenty swings are timed together.

\section*{Thermal Properties of Common Materials}

Table~\ref{tab:thermal} groups the materials by type. The specific heat is given for the
solid and the liquid phase where both are common in the laboratory.

\begin{table}[h]
\centering
\caption{Specific heat and latent heat of fusion.}
\label{tab:thermal}
\begin{tabular}{llccc}
\toprule
\multirow{2}{*}{Type} & \multirow{2}{*}{Material} & \multicolumn{2}{c}{Specific heat (\si{\joule\per\kilogram\per\kelvin})} & \multirow{2}{*}{$L_f$ (\si{\kilo\joule\per\kilogram})} \\
\cmidrule(lr){3-4}
 & & Solid & Liquid & \\
\midrule
\multirow{3}{*}{Metal} & Aluminium & 900 & 1180 & 397 \\
 & Copper & 385 & 530 & 205 \\
 & Lead & 128 & 140 & 24.5 \\
\midrule
\multirow{2}{*}{Other} & Water & 2090 & 4186 & 334 \\
 & Ethanol & 2420 & 2440 & 108 \\
\bottomrule
\end{tabular}
\end{table}

\clearpage
\section*{Summary of Kinematics and Dynamics}

The quantities below appear in every chapter of the course. Each definition is written
with the symbols used in the lectures.

\begin{table}[h]
\centering
\caption{Definitions of mechanical quantities.}
\label{tab:defs}
\renewcommand{\arraystretch}{2.1}
\begin{tabular}{llcc}
\toprule
Quantity & Symbol & Definition & SI unit \\
\midrule
Average velocity & $\bar{v}$ & $\displaystyle \bar{v} = \frac{\Delta x}{\Delta t}$ & \si{\metre\per\second} \\
Acceleration & $a$ & $\displaystyle a = \frac{dv}{dt}$ & \si{\metre\per\second\squared} \\
Momentum & $\vec{p}$ & $\vec{p} = m\vec{v}$ & \si{\kilogram\metre\per\second} \\
Kinetic energy & $K$ & $K = \tfrac{1}{2}mv^2$ & \si{\joule} \\
Work & $W$ & $W = \int \vec{F}\cdot d\vec{r}$ & \si{\joule} \\
\bottomrule
\end{tabular}
\end{table}

A quick check of units is the fastest way to catch an algebra mistake: the kinetic
energy has units of $\si{\kilogram}\,(\si{\metre\per\second})^2$, which is a joule.

\begin{table}[h]
\centering
\caption{Frequency of quiz scores.}
\label{tab:freq}
\begin{tabular}{|c|c|c|c|}
\hline
Score & Frequency & Relative frequency & Cumulative relative frequency \\
\hline
5 & 2 & 0.10 & 0.10 \\
\hline
6 & 3 & 0.15 & 0.25 \\
\hline
7 & 6 & 0.30 & 0.55 \\
\hline
8 & 5 & 0.25 & 0.80 \\
\hline
9 & 3 & 0.15 & 0.95 \\
\hline
10 & 1 & 0.05 & 1.00 \\
\hline
\end{tabular}
\end{table}

\clearpage
\section*{Derivatives by Finite Differences}

The table below gives values of $f(x) = x^3 - 2x$ and its derivative at evenly spaced
points. It is set without any rules, as many instructors type tables in a hurry.

\begin{table}[h]
\centering
\caption{Values of $f$ and $f'$.}
\label{tab:values}
\begin{tabular}{rrrr}
$x$ & $f(x)$ & $f'(x)$ & Forward difference \\[4pt]
$-2$ & $-4$ & $10$ & $5.0$ \\
$-1$ & $1$ & $1$ & $-1.0$ \\
$0$ & $0$ & $-2$ & $-1.0$ \\
$1$ & $-1$ & $1$ & $5.0$ \\
$2$ & $4$ & $10$ & $17.0$ \\
\end{tabular}
\end{table}

The forward difference with step $h = 1$ overestimates the derivative where the graph
curves upward, because
\begin{equation}
  \frac{f(x+h) - f(x)}{h} = f'(x) + \frac{h}{2} f''(\xi).
\end{equation}

\begin{table}[h]
\centering
\caption{Surface gravity of the planets.}
\label{tab:planets}
\begin{tabular}{>{\bfseries}lccc}
\toprule
\textbf{Planet} & \textbf{Mass ($10^{24}$ kg)} & \textbf{Radius (km)} & \textbf{$g$ (m/s$^2$)} \\
\midrule
Mercury & 0.330 & 2440 & 3.7 \\
Venus & 4.87 & 6052 & 8.9 \\
Earth & 5.97 & 6371 & 9.8 \\
Mars & 0.642 & 3390 & 3.7 \\
\bottomrule
\end{tabular}
\end{table}

\end{document}
